\documentclass[10pt,a4paper]{amsart}
\usepackage{amsmath,amssymb,amsthm,graphicx}
% For commutative diagrams we use Xy-pic
\usepackage[all]{xy}

% We use 2cell for 2-commutative diagrams.
\xyoption{2cell}
% Theorem environments.
%
\theoremstyle{plain}
\newtheorem{theorem}[subsection]{Theorem}
\newtheorem{proposition}[subsection]{Proposition}
\newtheorem{lemma}[subsection]{Lemma}

\theoremstyle{definition}
\newtheorem{definition}[subsection]{Definition}
\newtheorem{example}[subsection]{Example}
\newtheorem{exercise}[subsection]{Exercise}
\newtheorem{situation}[subsection]{Situation}

\theoremstyle{remark}
\newtheorem{remark}[subsection]{Remark}
\newtheorem{remarks}[subsection]{Remarks}

\numberwithin{equation}{subsection}

% Macros
%
\def\lim{\mathop{\mathrm{lim}}\nolimits}
\def\colim{\mathop{\mathrm{colim}}\nolimits}
\def\Spec{\mathop{\mathrm{Spec}}}
\def\Hom{\mathop{\mathrm{Hom}}\nolimits}
\def\Ext{\mathop{\mathrm{Ext}}\nolimits}
\def\SheafHom{\mathop{\mathcal{H}\!\mathit{om}}\nolimits}
\def\SheafExt{\mathop{\mathcal{E}\!\mathit{xt}}\nolimits}
\def\Sch{\mathit{Sch}}
\def\Mor{\mathop{\mathrm{Mor}}\nolimits}
\def\Ob{\mathop{\mathrm{Ob}}\nolimits}
\def\Sh{\mathop{\mathit{Sh}}\nolimits}
\def\NL{\mathop{N\!L}\nolimits}
\def\CH{\mathop{\mathrm{CH}}\nolimits}
\def\proetale{{pro\text{-}\acute{e}tale}}
\def\etale{{\acute{e}tale}}
\def\QCoh{\mathit{QCoh}}
\def\Ker{\mathop{\mathrm{Ker}}}
\def\Im{\mathop{\mathrm{Im}}}
\def\Coker{\mathop{\mathrm{Coker}}}
\def\Coim{\mathop{\mathrm{Coim}}}

% Boxtimes
%
\DeclareMathSymbol{\boxtimes}{\mathbin}{AMSa}{"02}

%
% Macros for moduli stacks/spaces
%
\def\QCohstack{\mathcal{QC}\!\mathit{oh}}
\def\Cohstack{\mathcal{C}\!\mathit{oh}}
\def\Spacesstack{\mathcal{S}\!\mathit{paces}}
\def\Quotfunctor{\mathrm{Quot}}
\def\Hilbfunctor{\mathrm{Hilb}}
\def\Curvesstack{\mathcal{C}\!\mathit{urves}}
\def\Polarizedstack{\mathcal{P}\!\mathit{olarized}}
\def\Complexesstack{\mathcal{C}\!\mathit{omplexes}}
% \Pic is the operator that assigns to X its picard group, usage \Pic(X)
% \Picardstack_{X/B} denotes the Picard stack of X over B
% \Picardfunctor_{X/B} denotes the Picard functor of X over B
\def\Pic{\mathop{\mathrm{Pic}}\nolimits}
\def\Picardstack{\mathcal{P}\!\mathit{ic}}
\def\Picardfunctor{\mathrm{Pic}}
\def\Deformationcategory{\mathcal{D}\!\mathit{ef}}
\newenvironment{reference}{}{}
\newenvironment{slogan}{}{}
\newenvironment{history}{}{}
\usepackage{cleveref}
\crefname{originalproposition}{original proposition}{original propositions}
\Crefname{originalproposition}{Original proposition}{Original propositions}
\makeatletter
\newlabel{spaces-limits-proposition-there-is-a-scheme-finite-over}{{09YC}{source}{}{sourcepointer-1}{}}
\newlabel{spaces-limits-proposition-there-is-a-scheme-finite-over@cref}{{[originalproposition][0][]09YC}{[1][0][]source}{}{}{}}
\makeatother
\crefname{originallemma}{original lemma}{original lemmas}
\Crefname{originallemma}{Original lemma}{Original lemmas}
\makeatletter
\newlabel{lemma-flatten-stack}{{0CQ4}{source}{}{sourcepointer-2}{}}
\newlabel{lemma-flatten-stack@cref}{{[originallemma][0][]0CQ4}{[1][0][]source}{}{}{}}
\makeatother
\crefname{originallemma}{original lemma}{original lemmas}
\Crefname{originallemma}{Original lemma}{Original lemmas}
\makeatletter
\newlabel{stacks-morphisms-lemma-gerbe-fppf}{{06QI}{source}{}{sourcepointer-3}{}}
\newlabel{stacks-morphisms-lemma-gerbe-fppf@cref}{{[originallemma][0][]06QI}{[1][0][]source}{}{}{}}
\makeatother
\crefname{originallemma}{original lemma}{original lemmas}
\Crefname{originallemma}{Original lemma}{Original lemmas}
\makeatletter
\newlabel{stacks-morphisms-lemma-gerbe-isom-fppf}{{0CPR}{source}{}{sourcepointer-4}{}}
\newlabel{stacks-morphisms-lemma-gerbe-isom-fppf@cref}{{[originallemma][0][]0CPR}{[1][0][]source}{}{}{}}
\makeatother
\crefname{originallemma}{original lemma}{original lemmas}
\Crefname{originallemma}{Original lemma}{Original lemmas}
\makeatletter
\newlabel{stacks-morphisms-lemma-separated-implies-isom}{{0CPL}{source}{}{sourcepointer-5}{}}
\newlabel{stacks-morphisms-lemma-separated-implies-isom@cref}{{[originallemma][0][]0CPL}{[1][0][]source}{}{}{}}
\makeatother
\crefname{originallemma}{original lemma}{original lemmas}
\Crefname{originallemma}{Original lemma}{Original lemmas}
\makeatletter
\newlabel{spaces-more-morphisms-lemma-get-section-after-blowup}{{0CPI}{source}{}{sourcepointer-6}{}}
\newlabel{spaces-more-morphisms-lemma-get-section-after-blowup@cref}{{[originallemma][0][]0CPI}{[1][0][]source}{}{}{}}
\makeatother
\crefname{originallemma}{original lemma}{original lemmas}
\Crefname{originallemma}{Original lemma}{Original lemmas}
\makeatletter
\newlabel{spaces-morphisms-lemma-equality-of-morphisms}{{084N}{source}{}{sourcepointer-7}{}}
\newlabel{spaces-morphisms-lemma-equality-of-morphisms@cref}{{[originallemma][0][]084N}{[1][0][]source}{}{}{}}
\makeatother
\crefname{originallemma}{original lemma}{original lemmas}
\Crefname{originallemma}{Original lemma}{Original lemmas}
\makeatletter
\newlabel{stacks-morphisms-lemma-section-immersion}{{050J}{source}{}{sourcepointer-8}{}}
\newlabel{stacks-morphisms-lemma-section-immersion@cref}{{[originallemma][0][]050J}{[1][0][]source}{}{}{}}
\makeatother
\crefname{originalproposition}{original proposition}{original propositions}
\Crefname{originalproposition}{Original proposition}{Original propositions}
\makeatletter
\newlabel{stacks-morphisms-proposition-open-stratum}{{06RC}{source}{}{sourcepointer-9}{}}
\newlabel{stacks-morphisms-proposition-open-stratum@cref}{{[originalproposition][0][]06RC}{[1][0][]source}{}{}{}}
\makeatother
\crefname{originallemma}{original lemma}{original lemmas}
\Crefname{originallemma}{Original lemma}{Original lemmas}
\makeatletter
\newlabel{stacks-limits-lemma-separated-closed-in-finite-presentation}{{0CQ2}{source}{}{sourcepointer-10}{}}
\newlabel{stacks-limits-lemma-separated-closed-in-finite-presentation@cref}{{[originallemma][0][]0CQ2}{[1][0][]source}{}{}{}}
\makeatother
\crefname{originalproposition}{original proposition}{original propositions}
\Crefname{originalproposition}{Original proposition}{Original propositions}
\makeatletter
\newlabel{spaces-limits-proposition-approximate}{{07SU}{source}{}{sourcepointer-11}{}}
\newlabel{spaces-limits-proposition-approximate@cref}{{[originalproposition][0][]07SU}{[1][0][]source}{}{}{}}
\makeatother
\crefname{originallemma}{original lemma}{original lemmas}
\Crefname{originallemma}{Original lemma}{Original lemmas}
\makeatletter
\newlabel{stacks-limits-lemma-descend-a-stack-down}{{0CN4}{source}{}{sourcepointer-12}{}}
\newlabel{stacks-limits-lemma-descend-a-stack-down@cref}{{[originallemma][0][]0CN4}{[1][0][]source}{}{}{}}
\makeatother
\crefname{originallemma}{original lemma}{original lemmas}
\Crefname{originallemma}{Original lemma}{Original lemmas}
\makeatletter
\newlabel{stacks-limits-lemma-eventually-separated}{{0CPZ}{source}{}{sourcepointer-13}{}}
\newlabel{stacks-limits-lemma-eventually-separated@cref}{{[originallemma][0][]0CPZ}{[1][0][]source}{}{}{}}
\makeatother
\crefname{originaldefinition}{original definition}{original definitions}
\Crefname{originaldefinition}{Original definition}{Original definitions}
\makeatletter
\newlabel{stacks-properties-definition-reduced-induced-stack}{{050C}{source}{}{sourcepointer-14}{}}
\newlabel{stacks-properties-definition-reduced-induced-stack@cref}{{[originaldefinition][0][]050C}{[1][0][]source}{}{}{}}
\makeatother
\crefname{originallemma}{original lemma}{original lemmas}
\Crefname{originallemma}{Original lemma}{Original lemmas}
\makeatletter
\newlabel{stacks-morphisms-lemma-gerbe-over-iso-classes}{{06QD}{source}{}{sourcepointer-15}{}}
\newlabel{stacks-morphisms-lemma-gerbe-over-iso-classes@cref}{{[originallemma][0][]06QD}{[1][0][]source}{}{}{}}
\makeatother
\crefname{originallemma}{original lemma}{original lemmas}
\Crefname{originallemma}{Original lemma}{Original lemmas}
\makeatletter
\newlabel{stacks-morphisms-lemma-flat-finite-presentation-permanence}{{06Q9}{source}{}{sourcepointer-16}{}}
\newlabel{stacks-morphisms-lemma-flat-finite-presentation-permanence@cref}{{[originallemma][0][]06Q9}{[1][0][]source}{}{}{}}
\makeatother
\crefname{originallemma}{original lemma}{original lemmas}
\Crefname{originallemma}{Original lemma}{Original lemmas}
\makeatletter
\newlabel{stacks-morphisms-lemma-gerbe-bijection-points}{{06R9}{source}{}{sourcepointer-17}{}}
\newlabel{stacks-morphisms-lemma-gerbe-bijection-points@cref}{{[originallemma][0][]06R9}{[1][0][]source}{}{}{}}
\makeatother
\crefname{originallemma}{original lemma}{original lemmas}
\Crefname{originallemma}{Original lemma}{Original lemmas}
\makeatletter
\newlabel{stacks-morphisms-lemma-locally-closed-in-noetherian}{{0CPM}{source}{}{sourcepointer-18}{}}
\newlabel{stacks-morphisms-lemma-locally-closed-in-noetherian@cref}{{[originallemma][0][]0CPM}{[1][0][]source}{}{}{}}
\makeatother
\crefname{originallemma}{original lemma}{original lemmas}
\Crefname{originallemma}{Original lemma}{Original lemmas}
\makeatletter
\newlabel{stacks-morphisms-lemma-check-separated-on-ui-cover}{{0CPQ}{source}{}{sourcepointer-19}{}}
\newlabel{stacks-morphisms-lemma-check-separated-on-ui-cover@cref}{{[originallemma][0][]0CPQ}{[1][0][]source}{}{}{}}
\makeatother
\crefname{originallemma}{original lemma}{original lemmas}
\Crefname{originallemma}{Original lemma}{Original lemmas}
\makeatletter
\newlabel{spaces-limits-lemma-embedding-into-affine-over-qs}{{088N}{source}{}{sourcepointer-20}{}}
\newlabel{spaces-limits-lemma-embedding-into-affine-over-qs@cref}{{[originallemma][0][]088N}{[1][0][]source}{}{}{}}
\makeatother
\crefname{originallemma}{original lemma}{original lemmas}
\Crefname{originallemma}{Original lemma}{Original lemmas}
\makeatletter
\newlabel{stacks-morphisms-lemma-factor-factor}{{0CPV}{source}{}{sourcepointer-21}{}}
\newlabel{stacks-morphisms-lemma-factor-factor@cref}{{[originallemma][0][]0CPV}{[1][0][]source}{}{}{}}
\makeatother
\crefname{originallemma}{original lemma}{original lemmas}
\Crefname{originallemma}{Original lemma}{Original lemmas}
\makeatletter
\newlabel{stacks-morphisms-lemma-scheme-theoretic-image-of-partial-section}{{0CPW}{source}{}{sourcepointer-22}{}}
\newlabel{stacks-morphisms-lemma-scheme-theoretic-image-of-partial-section@cref}{{[originallemma][0][]0CPW}{[1][0][]source}{}{}{}}
\makeatother
\crefname{originallemma}{original lemma}{original lemmas}
\Crefname{originallemma}{Original lemma}{Original lemmas}
\makeatletter
\newlabel{stacks-morphisms-lemma-existence-plus-flat-base-change}{{0CMK}{source}{}{sourcepointer-23}{}}
\newlabel{stacks-morphisms-lemma-existence-plus-flat-base-change@cref}{{[originallemma][0][]0CMK}{[1][0][]source}{}{}{}}
\makeatother
\crefname{originallemma}{original lemma}{original lemmas}
\Crefname{originallemma}{Original lemma}{Original lemmas}
\makeatletter
\newlabel{stacks-morphisms-lemma-universally-closed-permanence}{{0CPT}{source}{}{sourcepointer-24}{}}
\newlabel{stacks-morphisms-lemma-universally-closed-permanence@cref}{{[originallemma][0][]0CPT}{[1][0][]source}{}{}{}}
\makeatother
\crefname{originallemma}{original lemma}{original lemmas}
\Crefname{originallemma}{Original lemma}{Original lemmas}
\makeatletter
\newlabel{stacks-morphisms-lemma-topology-scheme-theoretic-image}{{0CML}{source}{}{sourcepointer-25}{}}
\newlabel{stacks-morphisms-lemma-topology-scheme-theoretic-image@cref}{{[originallemma][0][]0CML}{[1][0][]source}{}{}{}}
\makeatother
\crefname{originallemma}{original lemma}{original lemmas}
\Crefname{originallemma}{Original lemma}{Original lemmas}
\makeatletter
\newlabel{more-morphisms-lemma-dominate-fppf-finite}{{0CP0}{source}{}{sourcepointer-26}{}}
\newlabel{more-morphisms-lemma-dominate-fppf-finite@cref}{{[originallemma][0][]0CP0}{[1][0][]source}{}{}{}}
\makeatother
\crefname{originallemma}{original lemma}{original lemmas}
\Crefname{originallemma}{Original lemma}{Original lemmas}
\makeatletter
\newlabel{more-morphisms-lemma-extend-finite-surjective-morphisms}{{0CNY}{source}{}{sourcepointer-27}{}}
\newlabel{more-morphisms-lemma-extend-finite-surjective-morphisms@cref}{{[originallemma][0][]0CNY}{[1][0][]source}{}{}{}}
\makeatother
\title{Gabber–Chow proper covers and compactifications for separated finite-type algebraic stacks}
\author{The Stacks Project Authors}
\date{stacks-project:stacks-more-morphisms.tex:local-snapshot-20261011; GNU FDL 1.2 or later}
\begin{document}
\maketitle
\noindent\textit{Editorial note.} The source attributes the result to Ofer Gabber via Olsson. All named hypotheses are retained, especially separatedness; no assertion about arbitrary nonseparated classifying stacks is made.\par\medskip
\noindent\textit{Editorial note.} The previously prepared089J starts with an algebraic space and yields a quasi-projective scheme cover. This target starts with a stack and constructs an algebraic-space cover/compactification.\par\medskip
\noindent\textit{Allowed standard prerequisites (editorial).}\begin{itemize}
\item Explicit prerequisite boundary. Algebraic stacks and gerbes, finite-presentation approximation, reduced substacks, schematic images, properness and flat descent. The source-local make-section lemma is retained rather than imported as the desired cover theorem.
\end{itemize}\medskip
\ifdefined\hypertarget\hypertarget{sourcepointer-1}{}\fi
\noindent\textit{Source reference.} Original proposition 09YC: focused-seeds/AGP/sources/stacks-snapshot-20261011/tags.txt, source lines 12560--12560. Stable Stacks tag for an expressly allowed foundational or notational input; the counted construction remains the two concrete core obligations. The retained local core is not replaced by this pointer. See standard\_prerequisites.\par\medskip
\ifdefined\hypertarget\hypertarget{sourcepointer-2}{}\fi
\noindent\textit{Source reference.} Original lemma 0CQ4: focused-seeds/AGP/sources/stacks-snapshot-20261011/tags.txt, source lines 15949--15949. Stable Stacks tag for an expressly allowed foundational or notational input; the counted construction remains the two concrete core obligations. The retained local core is not replaced by this pointer. See standard\_prerequisites.\par\medskip
\ifdefined\hypertarget\hypertarget{sourcepointer-3}{}\fi
\noindent\textit{Source reference.} Original lemma 06QI: focused-seeds/AGP/sources/stacks-snapshot-20261011/tags.txt, source lines 8607--8607. Stable Stacks tag for an expressly allowed foundational or notational input; the counted construction remains the two concrete core obligations. The retained local core is not replaced by this pointer. See standard\_prerequisites.\par\medskip
\ifdefined\hypertarget\hypertarget{sourcepointer-4}{}\fi
\noindent\textit{Source reference.} Original lemma 0CPR: focused-seeds/AGP/sources/stacks-snapshot-20261011/tags.txt, source lines 15936--15936. Stable Stacks tag for an expressly allowed foundational or notational input; the counted construction remains the two concrete core obligations. The retained local core is not replaced by this pointer. See standard\_prerequisites.\par\medskip
\ifdefined\hypertarget\hypertarget{sourcepointer-5}{}\fi
\noindent\textit{Source reference.} Original lemma 0CPL: focused-seeds/AGP/sources/stacks-snapshot-20261011/tags.txt, source lines 15931--15931. Stable Stacks tag for an expressly allowed foundational or notational input; the counted construction remains the two concrete core obligations. The retained local core is not replaced by this pointer. See standard\_prerequisites.\par\medskip
\ifdefined\hypertarget\hypertarget{sourcepointer-6}{}\fi
\noindent\textit{Source reference.} Original lemma 0CPI: focused-seeds/AGP/sources/stacks-snapshot-20261011/tags.txt, source lines 15928--15928. Stable Stacks tag for an expressly allowed foundational or notational input; the counted construction remains the two concrete core obligations. The retained local core is not replaced by this pointer. See standard\_prerequisites.\par\medskip
\ifdefined\hypertarget\hypertarget{sourcepointer-7}{}\fi
\noindent\textit{Source reference.} Original lemma 084N: focused-seeds/AGP/sources/stacks-snapshot-20261011/tags.txt, source lines 10328--10328. Stable Stacks tag for an expressly allowed foundational or notational input; the counted construction remains the two concrete core obligations. The retained local core is not replaced by this pointer. See standard\_prerequisites.\par\medskip
\ifdefined\hypertarget\hypertarget{sourcepointer-8}{}\fi
\noindent\textit{Source reference.} Original lemma 050J: focused-seeds/AGP/sources/stacks-snapshot-20261011/tags.txt, source lines 6502--6502. Stable Stacks tag for an expressly allowed foundational or notational input; the counted construction remains the two concrete core obligations. The retained local core is not replaced by this pointer. See standard\_prerequisites.\par\medskip
\ifdefined\hypertarget\hypertarget{sourcepointer-9}{}\fi
\noindent\textit{Source reference.} Original proposition 06RC: focused-seeds/AGP/sources/stacks-snapshot-20261011/tags.txt, source lines 8636--8636. Stable Stacks tag for an expressly allowed foundational or notational input; the counted construction remains the two concrete core obligations. The retained local core is not replaced by this pointer. See standard\_prerequisites.\par\medskip
\ifdefined\hypertarget\hypertarget{sourcepointer-10}{}\fi
\noindent\textit{Source reference.} Original lemma 0CQ2: focused-seeds/AGP/sources/stacks-snapshot-20261011/tags.txt, source lines 15947--15947. Stable Stacks tag for an expressly allowed foundational or notational input; the counted construction remains the two concrete core obligations. The retained local core is not replaced by this pointer. See standard\_prerequisites.\par\medskip
\ifdefined\hypertarget\hypertarget{sourcepointer-11}{}\fi
\noindent\textit{Source reference.} Original proposition 07SU: focused-seeds/AGP/sources/stacks-snapshot-20261011/tags.txt, source lines 9914--9914. Stable Stacks tag for an expressly allowed foundational or notational input; the counted construction remains the two concrete core obligations. The retained local core is not replaced by this pointer. See standard\_prerequisites.\par\medskip
\ifdefined\hypertarget\hypertarget{sourcepointer-12}{}\fi
\noindent\textit{Source reference.} Original lemma 0CN4: focused-seeds/AGP/sources/stacks-snapshot-20261011/tags.txt, source lines 15879--15879. Stable Stacks tag for an expressly allowed foundational or notational input; the counted construction remains the two concrete core obligations. The retained local core is not replaced by this pointer. See standard\_prerequisites.\par\medskip
\ifdefined\hypertarget\hypertarget{sourcepointer-13}{}\fi
\noindent\textit{Source reference.} Original lemma 0CPZ: focused-seeds/AGP/sources/stacks-snapshot-20261011/tags.txt, source lines 15944--15944. Stable Stacks tag for an expressly allowed foundational or notational input; the counted construction remains the two concrete core obligations. The retained local core is not replaced by this pointer. See standard\_prerequisites.\par\medskip
\ifdefined\hypertarget\hypertarget{sourcepointer-14}{}\fi
\noindent\textit{Source reference.} Original definition 050C: focused-seeds/AGP/sources/stacks-snapshot-20261011/tags.txt, source lines 6495--6495. Stable Stacks tag for an expressly allowed foundational or notational input; the counted construction remains the two concrete core obligations. The retained local core is not replaced by this pointer. See standard\_prerequisites.\par\medskip
\ifdefined\hypertarget\hypertarget{sourcepointer-15}{}\fi
\noindent\textit{Source reference.} Original lemma 06QD: focused-seeds/AGP/sources/stacks-snapshot-20261011/tags.txt, source lines 8602--8602. Stable Stacks tag for an expressly allowed foundational or notational input; the counted construction remains the two concrete core obligations. The retained local core is not replaced by this pointer. See standard\_prerequisites.\par\medskip
\ifdefined\hypertarget\hypertarget{sourcepointer-16}{}\fi
\noindent\textit{Source reference.} Original lemma 06Q9: focused-seeds/AGP/sources/stacks-snapshot-20261011/tags.txt, source lines 8598--8598. Stable Stacks tag for an expressly allowed foundational or notational input; the counted construction remains the two concrete core obligations. The retained local core is not replaced by this pointer. See standard\_prerequisites.\par\medskip
\ifdefined\hypertarget\hypertarget{sourcepointer-17}{}\fi
\noindent\textit{Source reference.} Original lemma 06R9: focused-seeds/AGP/sources/stacks-snapshot-20261011/tags.txt, source lines 8633--8633. Stable Stacks tag for an expressly allowed foundational or notational input; the counted construction remains the two concrete core obligations. The retained local core is not replaced by this pointer. See standard\_prerequisites.\par\medskip
\ifdefined\hypertarget\hypertarget{sourcepointer-18}{}\fi
\noindent\textit{Source reference.} Original lemma 0CPM: focused-seeds/AGP/sources/stacks-snapshot-20261011/tags.txt, source lines 15932--15932. Stable Stacks tag for an expressly allowed foundational or notational input; the counted construction remains the two concrete core obligations. The retained local core is not replaced by this pointer. See standard\_prerequisites.\par\medskip
\ifdefined\hypertarget\hypertarget{sourcepointer-19}{}\fi
\noindent\textit{Source reference.} Original lemma 0CPQ: focused-seeds/AGP/sources/stacks-snapshot-20261011/tags.txt, source lines 15935--15935. Stable Stacks tag for an expressly allowed foundational or notational input; the counted construction remains the two concrete core obligations. The retained local core is not replaced by this pointer. See standard\_prerequisites.\par\medskip
\ifdefined\hypertarget\hypertarget{sourcepointer-20}{}\fi
\noindent\textit{Source reference.} Original lemma 088N: focused-seeds/AGP/sources/stacks-snapshot-20261011/tags.txt, source lines 10468--10468. Stable Stacks tag for an expressly allowed foundational or notational input; the counted construction remains the two concrete core obligations. The retained local core is not replaced by this pointer. See standard\_prerequisites.\par\medskip
\ifdefined\hypertarget\hypertarget{sourcepointer-21}{}\fi
\noindent\textit{Source reference.} Original lemma 0CPV: focused-seeds/AGP/sources/stacks-snapshot-20261011/tags.txt, source lines 15940--15940. Stable Stacks tag for an expressly allowed foundational or notational input; the counted construction remains the two concrete core obligations. The retained local core is not replaced by this pointer. See standard\_prerequisites.\par\medskip
\ifdefined\hypertarget\hypertarget{sourcepointer-22}{}\fi
\noindent\textit{Source reference.} Original lemma 0CPW: focused-seeds/AGP/sources/stacks-snapshot-20261011/tags.txt, source lines 15941--15941. Stable Stacks tag for an expressly allowed foundational or notational input; the counted construction remains the two concrete core obligations. The retained local core is not replaced by this pointer. See standard\_prerequisites.\par\medskip
\ifdefined\hypertarget\hypertarget{sourcepointer-23}{}\fi
\noindent\textit{Source reference.} Original lemma 0CMK: focused-seeds/AGP/sources/stacks-snapshot-20261011/tags.txt, source lines 15860--15860. Stable Stacks tag for an expressly allowed foundational or notational input; the counted construction remains the two concrete core obligations. The retained local core is not replaced by this pointer. See standard\_prerequisites.\par\medskip
\ifdefined\hypertarget\hypertarget{sourcepointer-24}{}\fi
\noindent\textit{Source reference.} Original lemma 0CPT: focused-seeds/AGP/sources/stacks-snapshot-20261011/tags.txt, source lines 15938--15938. Stable Stacks tag for an expressly allowed foundational or notational input; the counted construction remains the two concrete core obligations. The retained local core is not replaced by this pointer. See standard\_prerequisites.\par\medskip
\ifdefined\hypertarget\hypertarget{sourcepointer-25}{}\fi
\noindent\textit{Source reference.} Original lemma 0CML: focused-seeds/AGP/sources/stacks-snapshot-20261011/tags.txt, source lines 15861--15861. Stable Stacks tag for an expressly allowed foundational or notational input; the counted construction remains the two concrete core obligations. The retained local core is not replaced by this pointer. See standard\_prerequisites.\par\medskip
\ifdefined\hypertarget\hypertarget{sourcepointer-26}{}\fi
\noindent\textit{Source reference.} Original lemma 0CP0: focused-seeds/AGP/sources/stacks-snapshot-20261011/tags.txt, source lines 15910--15910. Stable Stacks tag for an expressly allowed foundational or notational input; the counted construction remains the two concrete core obligations. The retained local core is not replaced by this pointer. See standard\_prerequisites.\par\medskip
\ifdefined\hypertarget\hypertarget{sourcepointer-27}{}\fi
\noindent\textit{Source reference.} Original lemma 0CNY: focused-seeds/AGP/sources/stacks-snapshot-20261011/tags.txt, source lines 15908--15908. Stable Stacks tag for an expressly allowed foundational or notational input; the counted construction remains the two concrete core obligations. The retained local core is not replaced by this pointer. See standard\_prerequisites.\par\medskip
\begin{theorem}[Chow's lemma]
\label{theorem-chow-finite-type}
\begin{reference}
This is a result due to Ofer Gabber, see
\cite[Theorem 1.1]{olsson_proper}
\end{reference}
Let $f : \mathcal{X} \to Y$ be a morphism from an algebraic stack
to an algebraic space. Assume
\begin{enumerate}
\item $Y$ is quasi-compact and quasi-separated,
\item $f$ is separated of finite type.
\end{enumerate}
Then there exists a commutative diagram
$$
\xymatrix{
\mathcal{X} \ar[rd] & X \ar[l] \ar[d] \ar[r] & \overline{X} \ar[ld] \\
& Y
}
$$
where $X \to \mathcal{X}$ is proper surjective,
$X \to \overline{X}$ is an open immersion, and
$\overline{X} \to Y$ is proper morphism of algebraic spaces.
\end{theorem}

\section{Chow's lemma for algebraic stacks}
\label{section-chow}

\noindent
In this section we discuss Chow's lemma for algebraic stacks.

\begin{lemma}
\label{lemma-finite-cover-factor}
Let $Y$ be a quasi-compact and quasi-separated algebraic space.
Let $V \subset Y$ be a quasi-compact open. Let $f : \mathcal{X} \to V$
be surjective, flat, and locally of finite presentation.
Then there exists a finite surjective morphism $g : Y' \to Y$ such that
$V' = g^{-1}(V) \to Y$ factors Zariski locally through $f$.
\end{lemma}

\begin{proof}
We first prove this when $Y$ is a scheme.
We may choose a scheme $U$ and a surjective smooth morphism
$U \to \mathcal{X}$. Then $\{U \to V\}$ is an fppf covering of schemes.
By More on Morphisms, Lemma \ref{more-morphisms-lemma-dominate-fppf-finite}
there exists a finite surjective morphism
$V' \to V$ such that $V' \to V$ factors Zariski locally
through $U$. By
More on Morphisms, Lemma
\ref{more-morphisms-lemma-extend-finite-surjective-morphisms}
we can find a finite surjective morphism $Y' \to Y$
whose restriction to $V$ is $V' \to V$ as desired.

\medskip\noindent
If $Y$ is an algebraic space, then we see the lemma is true by
first doing a finite base change by a finite surjective morphism
$Y' \to Y$ where $Y'$ is a scheme. See
Limits of Spaces, Proposition
\ref{spaces-limits-proposition-there-is-a-scheme-finite-over}.
\end{proof}

\begin{lemma}
\label{lemma-make-section}
Let $f : \mathcal{X} \to Y$ be a morphism from an algebraic stack
to an algebraic space. Let $V \subset Y$ be an open subspace.
Assume
\begin{enumerate}
\item $f$ is separated and of finite type,
\item $Y$ is quasi-compact and quasi-separated,
\item $V$ is quasi-compact, and
\item $\mathcal{X}_V$ is a gerbe over $V$.
\end{enumerate}
Then there exists a commutative diagram
$$
\xymatrix{
\overline{Z} \ar[rd]_{\overline{g}} &
Z \ar[l]^j \ar[d]_g \ar[r]_h & \mathcal{X} \ar[ld]^f \\
& Y
}
$$
with $j$ an open immersion, $\overline{g}$ and $h$ proper,
and such that $|V|$ is contained in the image of $|g|$.
\end{lemma}

\begin{proof}
Suppose we have a commutative diagram
$$
\xymatrix{
\mathcal{X}' \ar[d]_{f'} \ar[r] & \mathcal{X} \ar[d]^f \\
Y' \ar[r] & Y
}
$$
and a quasi-compact open $V' \subset Y'$, such that
$Y' \to Y$ is a proper morphism of algebraic spaces,
$\mathcal{X}' \to \mathcal{X}$ is a proper morphism of algebraic stacks,
$V' \subset Y'$ maps surjectively onto $V$, and
$\mathcal{X}'_{V'}$ is a gerbe over $V'$.
Then it suffices to prove the lemma for the pair
$(f' : \mathcal{X}' \to Y', V')$. Some details omitted.

\medskip\noindent
Overall strategy of the proof. We will reduce
to the case where the image of $f$ is open and $f$
has a section over this open by repeatedly applying the
above remark. Each step is straightforward, but there are
quite a few of them which makes the proof a bit involved.

\medskip\noindent
Using Limits of Spaces, Proposition
\ref{spaces-limits-proposition-there-is-a-scheme-finite-over}
we reduce to the case where $Y$ is a scheme.
(Let $Y' \to Y$ be a finite surjective morphism where $Y'$ is
a scheme. Set $\mathcal{X}' = \mathcal{X}_{Y'}$ and apply
the initial remark of the proof.)

\medskip\noindent
Using Lemma \ref{lemma-flatten-stack}
(and Morphisms of Stacks, Lemma \ref{stacks-morphisms-lemma-gerbe-fppf}
to see that a gerbe is flat and locally of finite presentation)
we reduce to the case where $f$ is flat and of finite presentation.

\medskip\noindent
Since $f$ is flat and locally of finite
presentation, we see that the image of $|f|$ is an open $W \subset Y$.
Since $\mathcal{X}$ is quasi-compact (as $f$ is of finite type
and $Y$ is quasi-compact) we see that $W$ is quasi-compact.
By Lemma \ref{lemma-finite-cover-factor}
we can find a finite surjective morphism $g : Y' \to Y$
such that $g^{-1}(W) \to Y$ factors Zariski locally
through $\mathcal{X} \to Y$.
After replacing $Y$ by $Y'$ and $\mathcal{X}$ by
$\mathcal{X} \times_Y Y'$ we reduce to the situation
described in the next paragraph.

\medskip\noindent
Assume there exists $n \geq 0$, quasi-compact opens
$W_i \subset Y$, $i = 1, \ldots, n$, and
morphisms $x_i : W_i \to \mathcal{X}$ such that
(a) $f \circ x_i = \text{id}_{W_i}$,
(b) $W = \bigcup_{i = 1, \ldots, n} W_i$ contains $V$, and
(c) $W$ is the image of $|f|$.
We will use induction on $n$. The base case is $n = 0$: this
implies $V = \emptyset$ and in this case we can take
$\overline{Z} = \emptyset$.
If $n > 0$, then for $i = 1, \ldots, n$
consider the reduced closed subschemes
$Y_i$ with underlying topological space
$Y \setminus W_i$. Consider the finite morphism
$$
Y' = Y \amalg \coprod\nolimits_{i = 1, \ldots, n} Y_i \longrightarrow Y
$$
and the quasi-compact open
$$
V' = (W_1 \cap \ldots \cap W_n \cap V) \amalg
\coprod_{i = 1, \ldots, n} (V \cap Y_i).
$$
By the initial remark of the proof, if we can prove the lemma for the pairs
$$
(\mathcal{X} \to Y, W_1 \cap \ldots \cap W_n \cap V)
\quad\text{and}\quad
(\mathcal{X} \times_Y Y_i \to Y_i, V \cap Y_i),\quad
i = 1, \ldots, n
$$
then the result is true. Here we use the settheoretic equality
$V = (W_1 \cap \ldots \cap W_n \cap V) \cup
\bigcup\nolimits_{i = 1, \ldots n} (V \cap Y_i)$.
The induction hypothesis applies to the second type of
pairs above. Hence we reduce to the situation described in
the next paragraph.

\medskip\noindent
Assume there exists $n \geq 0$, quasi-compact opens
$W_i \subset Y$, $i = 1, \ldots, n$, and
morphisms $x_i : W_i \to \mathcal{X}$ such that
(a) $f \circ x_i = \text{id}_{W_i}$,
(b) $W = \bigcup_{i = 1, \ldots, n} W_i$ contains $V$,
(c) $W$ is the image of $|f|$, and
(d) $V \subset W_1 \cap \ldots \cap W_n$.
The morphisms
$$
T_{ij} = \mathit{Isom}_\mathcal{X}(x_i|_{W_i \cap W_j \cap V},
x_j|_{W_i \cap W_j \cap V}) \longrightarrow W_i \cap W_j \cap V
$$
are surjective, flat, and locally of finite presentation
(Morphisms of Stacks, Lemma \ref{stacks-morphisms-lemma-gerbe-isom-fppf}).
We apply Lemma \ref{lemma-finite-cover-factor}
to each quasi-compact open $W_i \cap W_j \cap V$ and the morphisms
$T_{ij} \to W_i \cap W_j \cap V$ to get finite surjective morphisms
$Y'_{ij} \to Y$. After replacing $Y$ by the fibre product of all
of the $Y'_{ij}$ over $Y$ we reduce to the situation described
in the next paragraph.

\medskip\noindent
Assume there exists $n \geq 0$, quasi-compact opens
$W_i \subset Y$, $i = 1, \ldots, n$, and
morphisms $x_i : W_i \to \mathcal{X}$ such that
(a) $f \circ x_i = \text{id}_{W_i}$,
(b) $W = \bigcup_{i = 1, \ldots, n} W_i$ contains $V$,
(c) $W$ is the image of $|f|$,
(d) $V \subset W_1 \cap \ldots \cap W_n$, and
(e) $x_i$ and $x_j$ are Zariski locally isomorphic over $W_i \cap W_j \cap V$.
Let $y \in V$ be arbitrary.
Suppose that we can find a quasi-compact open neighbourhood
$y \in V_y \subset V$ such that the lemma is true for
the pair $(\mathcal{X} \to Y, V_y)$, say with solution
$\overline{Z}_y, Z_y, \overline{g}_y, g_y, h_y$.
Since $V$ is quasi-compact, we can find a finite number
$y_1, \ldots, y_m$ such that $V = V_{y_1} \cup \ldots \cup V_{y_m}$.
Then it follows that setting
$$
\overline{Z} = \coprod \overline{Z}_{y_j},\quad
Z = \coprod Z_{y_j},\quad
\overline{g} = \coprod \overline{g}_{y_j},\quad
g = \coprod g_{y_j},\quad
h = \coprod h_{y_j}
$$
is a solution for the lemma. Given $y$ by condition (e)
we can choose a quasi-compact open neighbourhood $y \in V_y \subset V$
and isomorphisms $\varphi_i : x_1|_{V_y} \to x_i|_{V_y}$ for
$i = 2, \ldots, n$. Set $\varphi_{ij} = \varphi_j \circ \varphi_i^{-1}$.
This leads us to the situation described in the next paragraph.

\medskip\noindent
Assume there exists $n \geq 0$, quasi-compact opens
$W_i \subset Y$, $i = 1, \ldots, n$, and
morphisms $x_i : W_i \to \mathcal{X}$ such that
(a) $f \circ x_i = \text{id}_{W_i}$,
(b) $W = \bigcup_{i = 1, \ldots, n} W_i$ contains $V$,
(c) $W$ is the image of $|f|$,
(d) $V \subset W_1 \cap \ldots \cap W_n$, and
(f) there are isomorphisms $\varphi_{ij} : x_i|_V \to x_j|_V$
satisfying $\varphi_{jk} \circ \varphi_{ij} = \varphi_{ik}$.
The morphisms
$$
I_{ij} = \mathit{Isom}_\mathcal{X}(x_i|_{W_i \cap W_j},
x_j|_{W_i \cap W_j}) \longrightarrow W_i \cap W_j
$$
are proper because $f$ is separated
(Morphisms of Stacks, Lemma
\ref{stacks-morphisms-lemma-separated-implies-isom}).
Observe that $\varphi_{ij}$ defines a section $V \to I_{ij}$
of $I_{ij} \to W_i \cap W_j$ over $V$.
By More on Morphisms of Spaces, Lemma
\ref{spaces-more-morphisms-lemma-get-section-after-blowup}
we can find $V$-admissible blowups
$p_{ij} : Y_{ij} \to Y$ such that $s_{ij}$
extends to $p_{ij}^{-1}(W_i \cap W_j)$.
After replacing $Y$ by the fibre product of all the $Y_{ij}$
over $Y$ we get to the situation described in the next paragraph.

\medskip\noindent
Assume there exists $n \geq 0$, quasi-compact opens
$W_i \subset Y$, $i = 1, \ldots, n$, and
morphisms $x_i : W_i \to \mathcal{X}$ such that
(a) $f \circ x_i = \text{id}_{W_i}$,
(b) $W = \bigcup_{i = 1, \ldots, n} W_i$ contains $V$,
(c) $W$ is the image of $|f|$,
(d) $V \subset W_1 \cap \ldots \cap W_n$, and
(g) there are isomorphisms
$\varphi_{ij} : x_i|_{W_i \cap W_j} \to x_j|_{W_i \cap W_j}$
satisfying
$$
\varphi_{jk}|_V \circ \varphi_{ij}|_V = \varphi_{ik}|_V.
$$
After replacing $Y$ by another $V$-admissible blowup if necessary
we may assume that $V$ is dense and scheme theoretically dense
in $Y$ and hence in any open subspace of $Y$ containing $V$.
After such a replacement we conclude that
$$
\varphi_{jk}|_{W_i \cap W_j \cap W_k} \circ
\varphi_{ij}|_{W_i \cap W_j \cap W_k} =
\varphi_{ik}|_{W_i \cap W_j \cap W_k}
$$
by appealing to Morphisms of Spaces, Lemma
\ref{spaces-morphisms-lemma-equality-of-morphisms}
and the fact that $I_{ik} \to W_i \cap W_j$ is proper
(hence separated).
Of course this means that $(x_i, \varphi_{ij})$
is a desent datum and we obtain a morphism
$x : W \to \mathcal{X}$ agreeing with $x_i$ over $W_i$
because $\mathcal{X}$ is a stack.
Since $x$ is a section of the separated morphism
$\mathcal{X} \to W$ we see that $x$ is proper
(Morphisms of Stacks, Lemma \ref{stacks-morphisms-lemma-section-immersion}).
Thus the lemma now holds with $\overline{Z} = Y$,
$Z = W$, $\overline{g} = \text{id}_Y$, $g = \text{id}_W$,
$h = x$.
\end{proof}

\begin{proof}
The rough idea is to use that $\mathcal{X}$ has a dense open
which is a gerbe (Morphisms of Stacks, Proposition
\ref{stacks-morphisms-proposition-open-stratum})
and appeal to Lemma \ref{lemma-make-section}.
The reason this does not work is that the open may not be
quasi-compact and one runs into technical problems. Thus
we first do a (standard) reduction to the Noetherian case.

\medskip\noindent
First we choose a closed immersion
$\mathcal{X} \to \mathcal{X}'$ where $\mathcal{X}'$
is an algebraic stack separated and of finite type over $Y$.
See Limits of Stacks, Lemma
\ref{stacks-limits-lemma-separated-closed-in-finite-presentation}.
Clearly it suffices to prove the theorem for
$\mathcal{X}'$, hence we may assume $\mathcal{X} \to Y$
is separated and of finite presentation.

\medskip\noindent
Assume $\mathcal{X} \to Y$ is separated and of finite presentation.
By Limits of Spaces, Proposition \ref{spaces-limits-proposition-approximate}
we can write $Y = \lim Y_i$ as the directed limit of a system
of Noetherian algebraic spaces with affine transition morphisms.
By Limits of Stacks, Lemma \ref{stacks-limits-lemma-descend-a-stack-down}
there is an $i$ and a morphism $\mathcal{X}_i \to Y_i$ of finite presentation
from an algebraic stack to $Y_i$ such that
$\mathcal{X} = Y \times_{Y_i} \mathcal{X}_i$.
After increasing $i$ we may assume that $\mathcal{X}_i \to Y_i$
is separated, see Limits of Stacks, Lemma
\ref{stacks-limits-lemma-eventually-separated}.
Then it suffices to prove the theorem
for $\mathcal{X}_i \to Y_i$. This reduces us to the case discussed
in the next paragraph.

\medskip\noindent
Assume $Y$ is Noetherian. We may replace $\mathcal{X}$ by its
reduction (Properties of Stacks, Definition
\ref{stacks-properties-definition-reduced-induced-stack}).
This reduces us to the case discussed
in the next paragraph.

\medskip\noindent
Assume $Y$ is Noetherian and $\mathcal{X}$ is reduced.
Since $\mathcal{X} \to Y$ is separated and $Y$ quasi-separated,
we see that $\mathcal{X}$ is quasi-separated as an algebraic stack.
Hence the inertia $\mathcal{I}_\mathcal{X} \to \mathcal{X}$
is quasi-compact. Thus by Morphisms of Stacks, Proposition
\ref{stacks-morphisms-proposition-open-stratum}
there exists a dense open substack $\mathcal{V} \subset \mathcal{X}$
which is a gerbe. Let $\mathcal{V} \to V$ be the morphism
which expresses $\mathcal{V}$ as a gerbe over the algebraic space $V$.
See
Morphisms of Stacks, Lemma \ref{stacks-morphisms-lemma-gerbe-over-iso-classes}
for a construction of $\mathcal{V} \to V$.
This construction in particular shows that the morphism
$\mathcal{V} \to Y$ factors as $\mathcal{V} \to V \to Y$.
Picture
$$
\xymatrix{
\mathcal{V} \ar[r] \ar[d] & \mathcal{X} \ar[d] \\
V \ar[r] & Y
}
$$
Since the morphism $\mathcal{V} \to V$ is surjective, flat, and
of finite presentation
(Morphisms of Stacks, Lemma \ref{stacks-morphisms-lemma-gerbe-fppf})
and since $\mathcal{V} \to Y$ is locally of finite presentation,
it follows that $V \to Y$ is locally of finite presentation
(Morphisms of Stacks, Lemma
\ref{stacks-morphisms-lemma-flat-finite-presentation-permanence}).
Note that $\mathcal{V} \to V$ is a universal homeomorphism
(Morphisms of Stacks, Lemma
\ref{stacks-morphisms-lemma-gerbe-bijection-points}).
Since $\mathcal{V}$ is quasi-compact (see
Morphisms of Stacks, Lemma
\ref{stacks-morphisms-lemma-locally-closed-in-noetherian})
we see that $V$ is quasi-compact.
Finally, since $\mathcal{V} \to Y$ is separated the same is true
for $V \to Y$ by
Morphisms of Stacks, Lemma
\ref{stacks-morphisms-lemma-check-separated-on-ui-cover}
applied to $\mathcal{V} \to V \to Y$
(whose assumptions are satisfied as we've already seen).

\medskip\noindent
All of the above means that the assumptions of
Limits of Spaces, Lemma \ref{spaces-limits-lemma-embedding-into-affine-over-qs}
apply to the morphism $V \to Y$. Thus we can find a dense open
subspace $V' \subset V$ and an immersion $V' \to \mathbf{P}^n_Y$
over $Y$. Clearly we may replace $V$ by $V'$ and $\mathcal{V}$
by the inverse image of $V'$ in $\mathcal{V}$ (recall that
$|\mathcal{V}| = |V|$ as we've seen above).
Thus we may assume we have a diagram
$$
\xymatrix{
\mathcal{V} \ar[rr] \ar[d] & & \mathcal{X} \ar[d] \\
V \ar[r] & \mathbf{P}^n_Y \ar[r] & Y
}
$$
where the arrow $V \to \mathbf{P}^n_Y$ is an immersion.
Let $\mathcal{X}'$ be the scheme theoretic image of the morphism
$$
j : \mathcal{V} \longrightarrow \mathbf{P}^n_Y \times_Y \mathcal{X}
$$
and let $Y'$ be the scheme theoretic image of the morphism
$V \to \mathbf{P}^n_Y$. We obtain a commutative diagram
$$
\xymatrix{
\mathcal{V} \ar[r] \ar[d] &
\mathcal{X}' \ar[r] \ar[d] &
\mathbf{P}^n_Y \times_Y \mathcal{X} \ar[d] \ar[r] &
\mathcal{X} \ar[d] \\
V \ar[r] &
Y' \ar[r] &
\mathbf{P}^n_Y \ar[r] &
Y
}
$$
(See Morphisms of Stacks, Lemma \ref{stacks-morphisms-lemma-factor-factor}).
We claim that $\mathcal{V} = V \times_{Y'} \mathcal{X}'$ and that
Lemma \ref{lemma-make-section} applies to the morphism $\mathcal{X}' \to Y'$
and the open subspace $V \subset Y'$. If the claim is true, then we obtain
$$
\xymatrix{
\overline{X} \ar[rd]_{\overline{g}} &
X \ar[l] \ar[d]_g \ar[r]_h & \mathcal{X}' \ar[ld]^f \\
& Y'
}
$$
with $X \to \overline{X}$ an open immersion, $\overline{g}$ and $h$ proper,
and such that $|V|$ is contained in the image of $|g|$.
Then the composition $X \to \mathcal{X}' \to \mathcal{X}$ is
proper (as a composition of proper morphisms) and its image
contains $|\mathcal{V}|$, hence this composition is surjective.
As well, $\overline{X} \to Y' \to Y$ is proper as a composition
of proper morphisms.

\medskip\noindent
The last step is to prove the claim.
Observe that $\mathcal{X}' \to Y'$ is separated and of finite type,
that $Y'$ is quasi-compact and quasi-separated, and that $V$ is quasi-compact
(we omit checking all the details completely).
Next, we observe that
$b : \mathcal{X}' \to \mathcal{X}$ is an isomorphism over
$\mathcal{V}$ by Morphisms of Stacks, Lemma
\ref{stacks-morphisms-lemma-scheme-theoretic-image-of-partial-section}.
In particular $\mathcal{V}$ is identified with an open substack
of $\mathcal{X}'$.
The morphism $j$ is quasi-compact
(source is quasi-compact and target is quasi-separated), so formation
of the scheme theoretic image of $j$ commutes with flat base change by
Morphisms of Stacks, Lemma
\ref{stacks-morphisms-lemma-existence-plus-flat-base-change}.
In particular we see that $V \times_{Y'} \mathcal{X}'$ is the
scheme theoretic image of
$\mathcal{V} \to V \times_{Y'} \mathcal{X}'$.
However, by Morphisms of Stacks, Lemma
\ref{stacks-morphisms-lemma-universally-closed-permanence}
the image of $|\mathcal{V}| \to |V \times_{Y'} \mathcal{X}'|$
is closed (use that $\mathcal{V} \to V$ is a universal homeomorphism
as we've seen above and hence is universally closed).
Also the image is dense (combine what we just said with
Morphisms of Stacks, Lemma
\ref{stacks-morphisms-lemma-topology-scheme-theoretic-image})
we conclude $|\mathcal{V}| = |V \times_{Y'} \mathcal{X}'|$.
Thus $\mathcal{V} \to V \times_{Y'} \mathcal{X}'$ is an
isomorphism and the proof of the claim is complete.
\end{proof}

\bibliographystyle{plain}
\bibliography{my}
\end{document}
